% GENERATED FILE. Edit canonical lessons, then run npm run generate:books.
% canonical-source-hash: fnv1a64:166e9d7559275112
% canonical-lessons: ML-C323-tannimattan, ML-C323-cira, ML-C323-kararru, ML-C323-vellarikka, ML-C323-malli

\chapter{Watermelon, Spinach, Carrot, Cucumber, Coriander}
\label{ch:ml-watermelon-spinach-carrot-cucumber-coriander}
\hlchaptermodality{323}

\begin{chapteropening}
\textbf{By the end of this chapter:} \emph{I can say a watermelon, spinach, a carrot, a cucumber and coriander.}

Complete the last lesson of chapter 323: I can say a watermelon, spinach, a carrot, a cucumber and coriander.
\end{chapteropening}

\section[\texorpdfstring{\ml{തണ്ണിമത്തൻ} (taṇṇimattan)}{തണ്ണിമത്തൻ (taṇṇimattan)}]{\ml{തണ്ണിമത്തൻ} (taṇṇimattan) — a watermelon}

\label{lesson:ML-C323-tannimattan}



\begin{quote}
Before the new one: say the Malayalam for tapioca, then the Malayalam for a chutney.
\end{quote}

\subsection*{You'll want to know: \ml{തണ്ണിമത്തൻ}}
\textbf{\ml{തണ്ണിമത്തൻ}} — \emph{taṇṇimattan} — \textquotedblleft{}a watermelon\textquotedblright{}.

\begin{grammarlens}[title={What you've built}]
One more word for this chapter.
\end{grammarlens}

\subsection*{Your turn}
\begin{itemize}
\raggedright
  \item \emph{Say it:} \emph{taṇṇimattan}
  \item \emph{Say it:} \emph{taṇṇimattan}, once more
  \item \emph{Say it:} say \emph{taṇṇimattan}
  \item \emph{Recall:} say the Malayalam for tapioca, then the Malayalam for a chutney, then say \emph{taṇṇimattan} again
\end{itemize}

\begin{tcolorbox}[breakable,colback=teal!4,colframe=teal!35!black,title={Before you move on}]
What does \ml{തണ്ണിമത്തൻ} mean? (\textquotedblleft{}A watermelon\textquotedblright{}.) Say it once more.
\end{tcolorbox}
\section[\texorpdfstring{\ml{ചീര} (cīra)}{ചീര (cīra)}]{\ml{ചീര} (cīra) — spinach, greens}

\label{lesson:ML-C323-cira}



\begin{quote}
Before the new one: say the Malayalam for a chutney, then the Malayalam for a watermelon.
\end{quote}

\subsection*{You'll want to know: \ml{ചീര}}
\textbf{\ml{ചീര}} — \emph{cīra} — \textquotedblleft{}spinach, greens\textquotedblright{}.

\begin{grammarlens}[title={What you've built}]
One more word for this chapter.
\end{grammarlens}

\subsection*{Your turn}
\begin{itemize}
\raggedright
  \item \emph{Say it:} \emph{cīra}
  \item \emph{Say it:} \emph{cīra}, once more
  \item \emph{Say it:} say \emph{cīra}
  \item \emph{Recall:} say the Malayalam for a chutney, then the Malayalam for a watermelon, then say \emph{cīra} again
\end{itemize}

\begin{tcolorbox}[breakable,colback=teal!4,colframe=teal!35!black,title={Before you move on}]
What does \ml{ചീര} mean? (\textquotedblleft{}Spinach, greens\textquotedblright{}.) Say it once more.
\end{tcolorbox}
\section[\texorpdfstring{\ml{കാരറ്റ്} (kāraṟṟŭ)}{കാരറ്റ് (kāraṟṟŭ)}]{\ml{കാരറ്റ്} (kāraṟṟŭ) — a carrot}

\label{lesson:ML-C323-kararru}



\begin{quote}
Before the new one: say the Malayalam for a watermelon, then the Malayalam for spinach.
\end{quote}

\subsection*{You'll want to know: \ml{കാരറ്റ്}}
\textbf{\ml{കാരറ്റ്}} — \emph{kāraṟṟŭ} — \textquotedblleft{}a carrot\textquotedblright{}.

\begin{grammarlens}[title={What you've built}]
One more word for this chapter.
\end{grammarlens}

\subsection*{Your turn}
\begin{itemize}
\raggedright
  \item \emph{Say it:} \emph{kāraṟṟŭ}
  \item \emph{Say it:} \emph{kāraṟṟŭ}, once more
  \item \emph{Say it:} say \emph{kāraṟṟŭ}
  \item \emph{Recall:} say the Malayalam for a watermelon, then the Malayalam for spinach, then say \emph{kāraṟṟŭ} again
\end{itemize}

\begin{tcolorbox}[breakable,colback=teal!4,colframe=teal!35!black,title={Before you move on}]
What does \ml{കാരറ്റ്} mean? (\textquotedblleft{}A carrot\textquotedblright{}.) Say it once more.
\end{tcolorbox}
\section[\texorpdfstring{\ml{വെള്ളരിക്ക} (veḷḷarikka)}{വെള്ളരിക്ക (veḷḷarikka)}]{\ml{വെള്ളരിക്ക} (veḷḷarikka) — a cucumber}

\label{lesson:ML-C323-vellarikka}



\begin{quote}
Before the new one: say the Malayalam for spinach, then the Malayalam for a carrot.
\end{quote}

\subsection*{You'll want to know: \ml{വെള്ളരിക്ക}}
\textbf{\ml{വെള്ളരിക്ക}} — \emph{veḷḷarikka} — \textquotedblleft{}a cucumber\textquotedblright{}.

\begin{grammarlens}[title={What you've built}]
One more word for this chapter.
\end{grammarlens}

\subsection*{Your turn}
\begin{itemize}
\raggedright
  \item \emph{Say it:} \emph{veḷḷarikka}
  \item \emph{Say it:} \emph{veḷḷarikka}, once more
  \item \emph{Say it:} say \emph{veḷḷarikka}
  \item \emph{Recall:} say the Malayalam for spinach, then the Malayalam for a carrot, then say \emph{veḷḷarikka} again
\end{itemize}

\begin{tcolorbox}[breakable,colback=teal!4,colframe=teal!35!black,title={Before you move on}]
What does \ml{വെള്ളരിക്ക} mean? (\textquotedblleft{}A cucumber\textquotedblright{}.) Say it once more.
\end{tcolorbox}
\section[\texorpdfstring{\ml{മല്ലി} (malli)}{മല്ലി (malli)}]{\ml{മല്ലി} (malli) — coriander}

\label{lesson:ML-C323-malli}



\begin{quote}
Before the new one: say the Malayalam for a carrot, then the Malayalam for a cucumber.
\end{quote}

\subsection*{You'll want to know: \ml{മല്ലി}}
\textbf{\ml{മല്ലി}} — \emph{malli} — \textquotedblleft{}coriander\textquotedblright{}.

\begin{grammarlens}[title={What you've built}]
One more word for this chapter.
\end{grammarlens}

\subsection*{Your turn}
\begin{itemize}
\raggedright
  \item \emph{Say it:} \emph{malli}
  \item \emph{Say it:} \emph{malli}, once more
  \item \emph{Say it:} say \emph{malli}
  \item \emph{Recall:} say the Malayalam for a carrot, then the Malayalam for a cucumber, then say \emph{malli} again
\end{itemize}

\begin{tcolorbox}[breakable,colback=teal!4,colframe=teal!35!black,title={Before you move on}]
What does \ml{മല്ലി} mean? (\textquotedblleft{}Coriander\textquotedblright{}.) Say it once more.
\end{tcolorbox}
